\emph{Preliminaries.} By the relations of a switch (Theorem \claimref{synthesis.set-and-clear-of-a-bit-generate-the-hyperbolic-plane}(i), (iv)), $e_{a}^{2} = 1$, $f_{a}^{2} = -1$, $e_{a}f_{a} = w_{a}v_{a} - v_{a}w_{a} = 1 - 2n_{a} = -f_{a}e_{a}$, and by the hypothesis on distinct switches the $2m$ elements $e_{a}, f_{a}$ pairwise anticommute. Also $e_{a}v_{a} = w_{a}v_{a}$ and $v_{a}e_{a} = v_{a}w_{a}$, so $e_{a}v_{a}e_{a} = w_{a}v_{a}w_{a} = w_{a}(1 - w_{a}v_{a}) = w_{a}$, and likewise $e_{a}w_{a}e_{a} = v_{a}$.

\emph{(i)} A product of $m$ pairwise anticommuting elements $g_{1}, \dots, g_{m}$ squares to $(-1)^{m(m-1)/2}g_{1}^{2}\cdots g_{m}^{2}$, so $u^{2} = \pm 1$ and $u^{-1} = \pm u$. Let $m$ be odd, $u = e_{1}\cdots e_{m}$. Then $e_{a}$ commutes with itself and anticommutes with the other $m - 1$ factors, an even number, so $u\,e_{a} = e_{a}u$; and $f_{a}$ anticommutes with all $m$ factors, so $u\,f_{a} = -f_{a}u$. Let $m$ be even, $u = f_{1}\cdots f_{m}$. Then $f_{a}$ anticommutes with $m - 1$ factors, an odd number, so $u\,f_{a} = -f_{a}u$, and $e_{a}$ anticommutes with all $m$, so $u\,e_{a} = e_{a}u$. In both cases $X(e_{a}) = e_{a}$ and $X(f_{a}) = -f_{a}$, that is $X(v_{a} + w_{a}) = v_{a} + w_{a}$ and $X(v_{a} - w_{a}) = w_{a} - v_{a}$. For the null vectors themselves: for $m$ odd write $u = \pm e_{a}u'$ with $u'$ the product of the other $e_{b}$, an even number of elements each anticommuting with $v_{a}$, so $u'v_{a} = v_{a}u'$ and $u\,v_{a}\,u^{-1} = e_{a}v_{a}e_{a}^{-1} = e_{a}v_{a}e_{a} = w_{a}$ by the preliminaries; likewise $X(w_{a}) = v_{a}$. For $m$ even write $u = \pm f_{a}u'$ with $u'$ the product of the other $f_{b}$, an odd number, so $u'v_{a}u'^{-1} = -v_{a}$; and $f_{a}v_{a} = -w_{a}v_{a}$, so $f_{a}v_{a}f_{a} = -w_{a}v_{a}(v_{a} - w_{a}) = w_{a}v_{a}w_{a} = w_{a}$ and $f_{a}v_{a}f_{a}^{-1} = -f_{a}v_{a}f_{a} = -w_{a}$; together $X(v_{a}) = w_{a}$, and likewise $X(w_{a}) = v_{a}$. $X^{2}$ is conjugation by $u^{2} = \pm 1$, the identity. A ring homomorphism of the subring generated by the switches is determined by its values on the generators.

\emph{(ii)} The first statement is in the proof of (i), and multiplicativity gives $X(e_{S}f_{S'}) = (-1)^{|S'|}e_{S}f_{S'}$. $u$ is $e_{S}f_{S'}$ with $S$ or $S'$ the whole index set. The integer span of the blades is the image of the integer Clifford algebra under $e_{a} \mapsto v_{a} + w_{a}$, $f_{a} \mapsto v_{a} - w_{a}$ (Theorem \claimref{synthesis.integer-clifford-algebras-contain-no-switch}(iv) for one switch), and $X$ maps each blade to itself up to sign. $X(\alpha v_{a} + \beta w_{a}) = \beta v_{a} + \alpha w_{a}$, of the same square $\alpha\beta$. $X$ maps generators to generators, so it preserves the parity of a product of them. If $\rho$ is an anti-automorphism fixing the $v_{a}, w_{a}$, then $X\rho$ and $\rho X$ are anti-automorphisms that agree on the generators, sending $v_{a}$ to $w_{a}$ and $w_{a}$ to $v_{a}$; so they agree on the subring generated.

\emph{(iii)} $X(n_{a}) = w_{a}v_{a} = 1 - n_{a}$, so $X(h_{a}) = 2n_{a} - 1 = -h_{a}$, and the statements on the number, on $h_{1}\cdots h_{m}$ and on the two idempotents follow by additivity and multiplicativity.

\emph{(iv)} If $w_{a}\psi = 0$ and $X(\psi) = \psi$ then $v_{a}\psi = X(w_{a})X(\psi) = X(w_{a}\psi) = 0$, so $\psi = (v_{a}w_{a} + w_{a}v_{a})\psi = 0$. If $X(q) = -q$ and $\omega \circ X = \omega$ then $\omega(q) = \omega(X(q)) = -\omega(q)$.

\emph{(v)} If $\omega \circ X = \omega$ then $\omega(n_{1}) = \omega(1 - n_{1})$, so $2\,\omega(n_{1}) = \omega(1) = 1$, which no integer satisfies; only additivity, the normalisation and $X(n_{1}) = 1 - n_{1}$ are used. An additive map from a $\Q$-vector space to $\Z$ is zero, since each value is divisible by every positive integer, so a $\Q$-algebra has no integer state. Without the two inequalities the position need not be definite: on the integer span for $m = 1$, $2\omega_{1} - \omega_{0}$ is an integer state with values $2$ and $-1$ on $n_{1}$ and $1 - n_{1}$. With the adjoint, $v_{a}v_{a}^{\dagger} = v_{a}w_{a} = n_{a}$ and $w_{a}w_{a}^{\dagger} = w_{a}v_{a} = 1 - n_{a}$, so the two inequalities give $0 \le \omega(n_{a}) \le 1$ and $\omega(n_{a}) \in \{0, 1\}$; $q(\omega)$ is then a sum of $m$ terms $\pm 1$, and $q(\omega \circ X) = \sum_{a}\omega(1 - 2n_{a}) = -q(\omega)$. For the last statement, in $\Cl_{\Q}(m,m)$ the $n_{a}$ are commuting idempotents with $\langle p_{\sigma}\rangle_{0} = 2^{-m}$ (each factor is $\tfrac12(1 \pm e_{a}f_{a})$), so $\omega_{\sigma}(1) = 1$; $\omega_{\sigma}$ is the diagonal matrix entry at the configuration $\sigma$ in the representation on the $m$-fold product of two-state modules, in which the monomials in the $v_{a}, w_{a}$ have entries $0, \pm 1$, so it is integer-valued on their integer span; $p_{\sigma}n_{a} = \sigma_{a}p_{\sigma}$ gives $\omega_{\sigma}(n_{a}) = \sigma_{a}$ and the two inequalities; and $X(p_{\sigma}) = p_{\bar\sigma}$ with the invariance of the scalar part under conjugation gives $\omega_{\sigma} \circ X = \omega_{\bar\sigma}$. The charges $2|\sigma| - m$ run over all integers of absolute value at most $m$ and of the parity of $m$.

\emph{(vi)} The scalar part is invariant under conjugation, $n_{a} = \tfrac12(1 - e_{a}f_{a})$ has scalar part $\tfrac12$, and a blade has scalar part $0$ or $1$.
